\documentclass[11pt,a4paper]{article}
\usepackage[margin=25mm,headheight=15pt]{geometry}
\usepackage[T1]{fontenc}
\usepackage[utf8]{inputenc}
\IfFileExists{libertinus.sty}{\usepackage{libertinus}}{\usepackage{lmodern}}
\usepackage{microtype,amsmath,amssymb,amsthm,mathtools}
\usepackage{booktabs,longtable,array,tabularx}
\usepackage{xcolor,listings,fancyhdr,enumitem,xurl}
\usepackage[hidelinks]{hyperref}
\hypersetup{bookmarksdepth=2,pdftitle={Positive-real asymptotic inversion},pdfauthor={ChatGPT, prepared for Vladimir Reshetnikov},pdfsubject={Power-logarithmic inverse expansions and Wolfram Language implementation}}
\definecolor{ink}{HTML}{30352D}
\definecolor{codebg}{HTML}{F5F5F0}
\lstdefinelanguage{Wolfram}{morekeywords={Module,Block,Table,Do,If,Which,Return,Get,Clear,Total,Expand,FullSimplify,Log,Sqrt,D,Factorial,Assumptions,True,False,Failure,TestReport,VerificationTest,OptionsPattern,OptionValue},sensitive=true,morecomment=[s]{(*}{*)},morestring=[b]"}
\lstset{language=Wolfram,basicstyle=\small\ttfamily,columns=fullflexible,keepspaces=true,breaklines=true,showstringspaces=false,frame=single,rulecolor=\color{black!20},backgroundcolor=\color{codebg},xleftmargin=3pt,xrightmargin=3pt,aboveskip=9pt,belowskip=9pt}
\pagestyle{fancy}\fancyhf{}\fancyhead[L]{\small Positive-real asymptotic inversion}\fancyhead[R]{\small RealInverseAsymptotics 1.0}\fancyfoot[C]{\thepage}
\setlength{\parskip}{0.35em}\setlength{\parindent}{1em}
\setlist{itemsep=2pt,topsep=4pt}
\allowdisplaybreaks
\numberwithin{equation}{section}
\newtheorem{theorem}{Theorem}[section]
\newtheorem{lemma}[theorem]{Lemma}
\newtheorem{proposition}[theorem]{Proposition}
\newtheorem{corollary}[theorem]{Corollary}
\theoremstyle{definition}\newtheorem{definition}[theorem]{Definition}
\theoremstyle{remark}\newtheorem{remark}[theorem]{Remark}
\newcommand{\R}{\mathbb R}\newcommand{\N}{\mathbb N}
\newcommand{\Eul}{\mathcal E}
\newcommand{\val}{\operatorname{val}}
\newcommand{\supp}{\operatorname{supp}}
\newcommand{\PL}{\mathcal A_{\Delta}}
\newcommand{\dd}{\mathop{}\!\mathrm d}
\newcommand{\bigO}{\mathcal O}
\newcommand{\smallo}{\mathrm o}
\newcommand{\fall}[2]{(#1)^{\underline{#2}}}
\newcommand{\rising}[2]{(#1)^{\overline{#2}}}
\DeclareRobustCommand{\code}[1]{{\small\ttfamily\urlstyle{same}\nolinkurl{#1}}}

\title{\vspace{-12mm}\textbf{Positive-real asymptotic inversion}\\[5pt]
\Large Power-logarithmic series, exact irrational exponents,\\
and a Wolfram Language implementation}
\author{Technical article and implementation\\Prepared for Vladimir Reshetnikov}
\date{7 September 2026 \quad\textbullet\quad Version 1.0.0}

\begin{document}
\maketitle
\begin{abstract}
We construct the positive-real inverse branch at zero of
\(f(x)=a x^\alpha(1+\sum_j x^{\delta_j}P_j(\log x))\), where
\(a,\alpha,\delta_j>0\) and the \(P_j\) are real polynomials.
A universal Euler--Lagrange formula produces every coefficient without a
closed-form inverse, rationalizing exponents, or treating logarithms as
constants under differentiation. A finite positive exponent semigroup makes
weighted truncation an exact finite computation. A multivariate analytic
implicit-function argument proves convergence for the finite input class;
this is stronger than a merely formal expansion. We prove a first-omitted
power theorem that preserves the logarithmic factor in the error term, and
explain how inverse residuals and separately known input errors are transported.
The two motivating examples, \(x+x^2(1+\log x)\) and \(x+x^{\sqrt2}\), receive
explicit formulas, branch analysis, and numerical checks. The accompanying
Wolfram Language package implements sparse coefficient generation, exact
exponent comparison, conservative remainder metadata, and an independent
binomial/logarithm composition checker. Its scope is deliberately a precise
power-logarithmic class, rather than an unspecified universal transseries
solver.
\end{abstract}

\noindent\textbf{Deliverables and execution status.}
The archive contains this article and its source, the package, 35 Wolfram
regression tests, runnable examples, and independent validation scripts.
The mathematical algorithms passed 58 independently executed SymPy/mpmath
checks. A Wolfram kernel was unavailable in this session; the package and its
MUnit suite have not been executed in a proprietary Wolfram runtime.
Section~\ref{sec:validation} separates these statements in detail.

\newpage
\setcounter{tocdepth}{1}
{\small\setlength{\parskip}{0pt}\tableofcontents}
\newpage

\section{The question is about a germ, not just an inverse symbol}
\label{sec:question}

The supplied question~\cite{question} asks for the real inverse of
\begin{equation}\label{eq:f1}
 f_1(x)=x+x^2(1+\log x),\qquad x>0,
\end{equation}
near the image point zero. A second example asks for the inverse of
\begin{equation}\label{eq:f2}
 f_2(x)=x+x^{\sqrt2},\qquad x\ge0.
\end{equation}
The archived evaluation of an unqualified \code{InverseFunction} produces
an expansion with a nonreal constant term in the first example. This does
not identify the desired germ. The specification needed here is
\begin{equation}\label{eq:branch-unit}
 f(g(y))=y,\qquad y\downarrow0,\qquad g(y)>0,\qquad g(y)\sim y.
\end{equation}
The target point alone does not specify a preimage or a branch. In particular,
putting a reality condition around a symbolic expression is not the same as
constructing the local inverse germ with a stated asymptotic leading term.
We analyze the archived output, not the behavior of an untested current
Wolfram release.

\subsection{The first example actually has a unique positive inverse globally}
Differentiation gives
\[
 f_1'(x)=1+x(3+2\log x).
\]
The function \(x(3+2\log x)\) has its global minimum at \(x=e^{-5/2}\),
where its value is \(-2e^{-5/2}\). Consequently
\begin{equation}\label{eq:global-slope}
 f_1'(x)\ge 1-2e^{-5/2}>0\qquad(x>0).
\end{equation}
Moreover \(f_1(x)\to0\) as \(x\downarrow0\), and \(f_1(x)\to\infty\) as
\(x\to\infty\). Thus it is a bijection of the positive half-line onto itself.
There is no ambiguity about its positive-real inverse. The difficulty is
representing its nonanalytic endpoint germ.

For comparison, a nonzero complex zero of the analytic continuation satisfies
\(x(1+\log x)=-1\). Writing \(w=1+\log x\) gives \(we^w=-e\), hence
\[
 x=\exp\!\bigl(W_k(-e)-1\bigr)=-\frac1{W_k(-e)}
\]
where the chosen logarithm and Lambert branch must be compatible. The
nonreal constant displayed in the archived question is of this form. An
ordinary analytic inverse expansion there belongs to a different preimage
of zero, not to the positive endpoint branch.

For \(f_2\), the derivative \(1+\sqrt2\,x^{\sqrt2-1}\) is positive for
\(x>0\). This also gives a unique positive inverse, globally.

\subsection{Two independent representation issues}
The logarithmic example does not have an ordinary Taylor expansion at zero.
For instance, the third derivative of \(x^2\log x\) is \(2/x\). The second
example introduces the irrational exponent gap \(\sqrt2-1\); no integer
ramification turns all its inverse exponents into integers.

The documented \code{SeriesData} representation uses one lattice of powers
\((n_{\min}+j)/\mathrm{den}\), and \code{InverseSeries} operates on suitably
structured objects of that kind~\cite{seriesdata,inverseseries}. This is not
the representation used below. A sparse list of exact real exponents with
polynomial-logarithmic coefficients expresses both examples directly.
Wolfram's \code{Asymptotic} documentation allows general inferred scales,
but that does not remove the need to specify the inverse germ or establish
the appropriate coefficient algebra~\cite{asymptotic}.

\subsection{How this relates to the supplied transseries material}
The ProveIt material~\cite{proveit} distinguishes support conditions, formal
summability, compositional inversion, logarithmic degree growth, and
analytic remainder claims. Its chapter on the Lagrange--B\"urmann formula at
infinity is particularly pertinent. We use the same separation, but work at
\(0^+\), allow finitely many exact irrational exponent generators, normalize
an arbitrary positive leading monomial, and prove an analytic realization
for the finite input class.

The broader problem belongs to the theory of transseries and multiseries.
Edgar's introduction and van der Hoeven's monograph provide background;
Salvy and Shackell specifically study algorithms for multiseries expansions
of inverse exp-log functions~\cite{edgar,hoeven,salvy}. The package here is a
specialized implementation with a narrower, explicit contract. Its proofs
below do not assume a black-box theorem saying that every imaginable
transseries can be substituted or inverted.

\section{The input class and the real branch}
\label{sec:class}

Fix real constants \(a>0\), \(\alpha>0\), and a finite collection of positive
numbers \(\delta_1,\ldots,\delta_r\). Consider
\begin{equation}\label{eq:class}
 f(x)=a x^\alpha\bigl(1+u(x)\bigr),\qquad
 u(x)=\sum_{j=1}^r x^{\delta_j}P_j(\log x),\qquad P_j\in\R[L].
\end{equation}
Combine equal \(\delta_j\) and discard zero polynomials. The coefficients
may depend on fixed real parameters; all assertions are then for parameter
values satisfying the stated hypotheses, not automatically uniform as a
parameter approaches a singular boundary.

Throughout, real powers mean \(x^q=\exp(q\log x)\) for \(x>0\). Set
\begin{equation}\label{eq:normalization}
 t=(y/a)^{1/\alpha},\qquad L=\log t,
 \qquad\Eul=t\frac{\dd}{\dd t}.
\end{equation}
The required branch is the one satisfying \(g(y)/t\to1\).
The symbol \(L\) is an independent polynomial coordinate during coefficient
arithmetic, but its differential rule is never forgotten:
\begin{equation}\label{eq:euler-block}
 \Eul\bigl[t^\mu Q(L)\bigr]
 =t^\mu\bigl(\mu+\partial_L\bigr)Q(L).
\end{equation}

\begin{theorem}[Existence and selection of the positive branch]\label{thm:branch}
For a nonzero finite perturbation in \eqref{eq:class}, let
\(\delta_* =\min_j\delta_j\) and \(d=\max_j\deg P_j\).
For sufficiently small \(\varepsilon>0\), \(f\) is positive and strictly
increasing on \((0,\varepsilon)\), extends continuously by \(f(0)=0\), and
has a unique positive inverse there. That inverse satisfies
\[
 g(y)\sim (y/a)^{1/\alpha},\qquad
 g'(y)\sim \frac{(y/a)^{1/\alpha}}{\alpha y}.
\]
\end{theorem}
\begin{proof}
Every fixed power of \(|\log x|\) is dominated by every positive power of
\(x^{-1}\). Therefore
\[
 u(x)=\bigO\!\left(x^{\delta_*}(1+|\log x|)^d\right)=\smallo(1),
 \qquad xu'(x)=\smallo(1).
\]
The derivative is
\[
 f'(x)=a x^{\alpha-1}
 \left[\alpha(1+u(x))+xu'(x)\right]
 \sim a\alpha x^{\alpha-1}>0.
\]
This proves local positivity and monotonicity. From
\(y=a g(y)^\alpha(1+u(g(y)))\) we obtain \(g(y)/t\to1\).
Finally \(g'=1/f'(g)\) yields the derivative equivalent. The argument does
not require a finite nonzero derivative of \(f\) at the endpoint; thus
\(\alpha<1\) and \(\alpha>1\) are both allowed.
\end{proof}

\section{The coefficient algebra and why truncation is finite}
\label{sec:algebra}

\subsection{A positive semigroup, not a rational exponent lattice}
Let
\begin{equation}\label{eq:semigroup}
 S=\left\{m_1\delta_1+\cdots+m_r\delta_r:
           (m_1,\ldots,m_r)\in\N^r\right\}.
\end{equation}
For every finite \(B\), only finitely many multiindices can satisfy
\(\sum m_j\delta_j\le B\), because \(m_j\le B/\delta_j\).
Thus \(S\) is locally finite. This remains true when some ratios
\(\delta_i/\delta_j\) are irrational. The additive \emph{group} they generate
may be dense; the nonnegative \emph{semigroup} relevant here is locally finite.
Confusing these two statements can lead to an incorrect objection to
finite coefficient extraction.

Define the formal algebra
\[
 \PL=\left\{\sum_{\mu\in S}t^\mu Q_\mu(L):Q_\mu\in\R[L]\right\},
 \qquad
 \val A=\min\{\mu:Q_\mu\ne0\}.
\]
Each bounded power interval contains only finitely many polynomial blocks.
Multiplication is the finite convolution
\[
 [t^\mu](AB)=\sum_{\lambda+\nu=\mu}A_\lambda(L)B_\nu(L).
\]
The Euler derivation preserves power valuation and does not increase
logarithmic degree. A fixed coefficient below a power cutoff is therefore
computed with finitely many exact polynomial operations.

\subsection{Composition near the identity}
For \(v\) of positive valuation, the expressions
\begin{equation}\label{eq:unit-functions}
 (1+v)^q=\sum_{n\ge0}\binom qn v^n,
 \qquad
 \log(1+v)=\sum_{n\ge1}\frac{(-1)^{n+1}}n v^n
\end{equation}
are formally summable. Substitution into a block uses
\begin{equation}\label{eq:block-substitution}
 (t(1+v))^\delta P\bigl(\log(t(1+v))\bigr)
 =t^\delta(1+v)^\delta P\bigl(L+\log(1+v)\bigr).
\end{equation}
This gives a complete rule for composition within the class. In particular,
it differentiates and shifts logarithms correctly; replacing \(\log t\) by
an unrelated constant before reversing a series would not implement this rule.

\begin{proposition}[Formal existence and uniqueness]\label{prop:formal}
There is a unique \(s\in1+\mathcal A_{\Delta,>0}\) satisfying
\begin{equation}\label{eq:normalized-inverse}
 s^\alpha\bigl(1+u(ts)\bigr)=1.
\end{equation}
Its support is contained in \(S\), and the formal inverse is \(g=t s\).
\end{proposition}
\begin{proof}
Rewrite the equation as
\(s=(1+u(ts))^{-1/\alpha}\), with the binomial branch having constant term
one. If \(s_1-s_2\) has valuation \(\eta>0\), formula
\eqref{eq:block-substitution} shows that
\(u(ts_1)-u(ts_2)\) has valuation at least \(\eta+\delta_*\).
Applying the unit power preserves that gain. Starting from \(s_0=1\),
iteration therefore stabilizes every bounded set of coefficients after
finitely many steps. The same valuation gain proves uniqueness: a nonzero
least-weight difference of two solutions would have to gain positive weight
while remaining equal to itself. All operations preserve the semigroup.
\end{proof}

This proposition also applies to locally finite infinite formal inputs
with a strictly positive least correction weight. It says nothing by itself
about convergence or the analytic size of a remainder. Those require the
separate arguments in Sections~\ref{sec:convergence} and~\ref{sec:truncation}.

\section{A universal Euler--Lagrange inversion formula}
\label{sec:formula}

\begin{theorem}[Euler--Lagrange formula]\label{thm:euler}
For \eqref{eq:class}, the normalized inverse is
\begin{equation}\label{eq:master}
 \boxed{\displaystyle
 g(y)=t\left[1+\sum_{n\ge1}
 \frac{(-1)^n}{n!\,\alpha^n}
 \prod_{j=1}^{n-1}(\Eul+1+j\alpha)\,u(t)^n\right].}
\end{equation}
An empty operator product is the identity. The formula is valid formally;
for the finite input class it also converges for sufficiently small positive
\(y\), as proved in Theorem~\ref{thm:convergence}.
\end{theorem}

\subsection{Derivation without assuming analyticity at zero}
Classical Lagrange inversion is an analytic local theorem at a regular
point~\cite{dlmf}. We apply it in an auxiliary perturbation parameter, around
a positive point, not at the singular endpoint.

Introduce \(\lambda\) and consider
\[
 x^\alpha(1+\lambda u(x))=t^\alpha.
\]
For fixed \(t>0\), put \(w=t^\alpha\), \(z=x^\alpha\), and
\[
 R(z)=z\,u(z^{1/\alpha}),\qquad H(z)=z^{1/\alpha}.
\]
All functions are analytic near the positive point \(z=w\), with consistent
local branches. The equation becomes \(z+\lambda R(z)=w\).
The parameter form of Lagrange--B\"urmann gives
\begin{equation}\label{eq:observable-lagrange}
 H(z(\lambda))=H(w)+\sum_{n\ge1}\frac{(-\lambda)^n}{n!}
 \left(\frac{\dd}{\dd w}\right)^{n-1}
       \left[H'(w)R(w)^n\right].
\end{equation}
For completeness, the coefficient identity follows by the residue formula
for \(H(z(\lambda))\), changing variables from the inverse to
\(z+\lambda R(z)-w\), and integrating the derivative of the logarithm once
by parts. At order \(\lambda^n\) the residue is
\[
 \frac{(-1)^n}{n}\operatorname*{Res}_{z=w}
 \frac{H'(z)R(z)^n}{(z-w)^n}
 =\frac{(-1)^n}{n!}
 \left(\frac{\dd}{\dd w}\right)^{n-1}[H'(w)R(w)^n].
\]
Equivalently, this is the usual formal-parameter reversion identity; no
convergence of a Taylor series in \(t\) about zero has been invoked.

Here
\[
 H'(w)R(w)^n=\frac1\alpha
 w^{n-1+1/\alpha}u(w^{1/\alpha})^n.
\]
Repeated differentiation, followed by \(w=t^\alpha\), yields
\begin{equation}\label{eq:euler-transform}
 \left(\frac{\dd}{\dd w}\right)^{n-1}
 \left[w^{n-1+1/\alpha}A(w^{1/\alpha})\right]
 =\frac{t}{\alpha^{n-1}}
   \prod_{j=1}^{n-1}(\Eul+1+j\alpha)A(t).
\end{equation}
For a monomial \(A(t)=t^\mu\), both sides have the same product of
\(n-1\) linear factors. The identity for general differentiable \(A\)
follows by iterating \(\dd/\dd w=(\alpha t^{\alpha-1})^{-1}\dd/\dd t\).
Substitution into \eqref{eq:observable-lagrange} gives \eqref{eq:master}.
Formal summability follows because its \(n\)-th correction has relative
valuation at least \(n\delta_*\). Setting \(\lambda=1\) analytically is
justified separately below.

\subsection{The familiar near-identity formula}
When \(a=\alpha=1\), write \(h(t)=t u(t)\), so that \(f(t)=t+h(t)\).
The Euler formula becomes
\begin{equation}\label{eq:unit-master}
 \boxed{\displaystyle
 g(y)=y+\sum_{n\ge1}\frac{(-1)^n}{n!}
             \left(\frac{\dd}{\dd y}\right)^{n-1}h(y)^n.}
\end{equation}
Indeed
\(\dd^{n-1}(t^n u^n)/\dd t^{n-1}
=t\prod_{j=2}^{n}(\Eul+j)u^n\).
This short formula already solves both examples. The package adds support
management, normalization, exact cutoffs, branch metadata, and independent
verification.

A compact Wolfram implementation for a fixed number of correction orders is
\begin{lstlisting}
Clear[UnitInverseJet];
UnitInverseJet[h_, x_Symbol, n_Integer] /; n >= 0 :=
  x + Total[Table[(-1)^k D[h^k, {x, k-1}]/k!, {k, 1, n}]];

Expand[UnitInverseJet[x^2 (1+Log[x]), x, 3]]
\end{lstlisting}
This expression does not attach an error term and does not enforce a power
cutoff. The distinction matters for a mixture of different exponent gaps.

\subsection{A coefficient formula on exact sparse blocks}
Write
\(u(t)^n=\sum_\mu t^\mu Q_{n,\mu}(L)\).
The contribution of its \((n,\mu)\) block is
\begin{equation}\label{eq:block-coefficient}
 t^{1+\mu}\frac{(-1)^n}{n!\alpha^n}
 \prod_{j=1}^{n-1}\left(\mu+1+j\alpha+\partial_L\right)Q_{n,\mu}(L).
\end{equation}
Only polynomial multiplication and polynomial differentiation are needed.
The operators commute and do not raise logarithmic degree.

For a multiindex \(m=(m_1,\ldots,m_r)\), put
\(n=|m|\), \(\mu=m\cdot\delta\), and \(m!=\prod m_j!\). Expanding
\(u^n\) by the multinomial theorem gives the nonrecursive version
\begin{equation}\label{eq:multiindex}
 g=t+\sum_{m\ne0}
 \frac{(-1)^{|m|}\,t^{1+m\cdot\delta}}{\alpha^{|m|}m!}
 \prod_{j=1}^{|m|-1}
 \left(m\cdot\delta+1+j\alpha+\partial_L\right)
 \prod_{i=1}^r P_i(L)^{m_i}.
\end{equation}
Different multiindices can give the same exponent. Their polynomial
coefficients must be added \emph{before} deciding whether a term vanishes.
This is the source of resonances and cancellations even for perfectly
ordinary integer exponents.

\subsection{Degree bounds and pure-power coefficients}
For a fixed multiindex, the logarithmic degree is at most
\(\sum m_i\deg P_i\). In fact, before cancellations between multiindices,
the degree of a nonzero polynomial in \eqref{eq:multiindex} is unchanged:
each factor has a strictly positive constant part
\(m\cdot\delta+1+j\alpha\).
In particular, a block of relative weight \(\mu\) has degree at most
\begin{equation}\label{eq:degree-bound}
 d\left\lfloor\frac\mu{\delta_*}\right\rfloor.
\end{equation}
This controls coefficient sizes, not analytic tail sizes by itself.

If every \(P_i=c_i\) is constant, the coefficient is especially simple:
\begin{equation}\label{eq:pure-multi}
 [t^{1+\mu}]g
 =\sum_{\substack{m\ne0\\m\cdot\delta=\mu}}
 \frac{(-1)^{|m|}}{\alpha^{|m|}m!}
 \prod_{j=1}^{|m|-1}(\mu+1+j\alpha)\prod_i c_i^{m_i}.
\end{equation}
The sum is finite. No common denominator for the \(\delta_i\) is needed.

\section{Analytic realization and convergence}
\label{sec:convergence}

\begin{theorem}[Convergence for finite power-logarithmic input]\label{thm:convergence}
For \eqref{eq:class} with finitely many polynomial-logarithmic terms, the
series \eqref{eq:master}, evaluated at \(t=(y/a)^{1/\alpha}\), converges for
all sufficiently small positive \(y\) to the inverse in
Theorem~\ref{thm:branch}. It can be regrouped by increasing exact power
weight. If all correction orders \(n<N\) are retained, the tail is
\begin{equation}\label{eq:degree-tail}
 \bigO\!\left(t^{1+N\delta_*}(1+|\log t|)^{dN}\right).
\end{equation}
\end{theorem}
\begin{proof}
Write \(g=t s\), so the defining equation is \eqref{eq:normalized-inverse}.
For each \(j\) and \(0\le k\le d_j=\deg P_j\), introduce independent small
variables
\[
 z_{j,k}=t^{\delta_j}L^k.
\]
The polynomial identity
\[
 P_j(L+\log s)=\sum_{k=0}^{d_j}\frac{L^k}{k!}P_j^{(k)}(\log s)
\]
turns the inverse equation into the finite-dimensional analytic equation
\begin{equation}\label{eq:analytic-lift}
 \Phi(s,z)=s^\alpha\left[
 1+\sum_{j=1}^r s^{\delta_j}
      \sum_{k=0}^{d_j}\frac{z_{j,k}}{k!}P_j^{(k)}(\log s)
 \right]-1=0.
\end{equation}
Use the branches of \(s^q\) and \(\log s\) analytic in a disk about \(s=1\).
At \((s,z)=(1,0)\), we have \(\Phi=0\) and \(\partial_s\Phi=\alpha\ne0\).
The analytic implicit-function theorem gives an analytic function \(s(z)\)
in a fixed polydisk about zero, with \(s(0)=1\).

Every \(z_{j,k}(t)\to0\) as \(t\downarrow0\). Thus the actual path enters
that polydisk and \(t s(z(t))\) is the positive inverse. Scaling all the
variables simultaneously by \(\lambda\) replaces \(u\) by \(\lambda u\).
The coefficient of \(\lambda^n\) in \(t s(\lambda z(t))\) is exactly the
\(n\)-th Euler--Lagrange correction. For small \(t\), its analytic radius in
\(\lambda\) exceeds one, so evaluation at \(\lambda=1\) is justified.

The multivariate Taylor series is absolutely convergent inside a smaller
polydisk. It may therefore be regrouped by total degree and by the locally
finite power weights. For a quantitative estimate, choose a closed smaller
polydisk on which \(|s|\) is bounded, and apply the one-variable Cauchy
estimate along \(\lambda z\). With constants independent of small \(t\),
\[
 \left|\sum_{n\ge N}[\lambda^n]\,t s(\lambda z(t))\right|
 \le C t\,q(t)^N,
 \qquad
 q(t)\le C_1 t^{\delta_*}(1+|L|)^d<\tfrac12.
\]
This proves \eqref{eq:degree-tail}.
\end{proof}

\begin{remark}[What is convergent here]
For the logarithmic example, the result is not an ordinary Taylor series in
\(y\). It is a convergent regrouping of a multivariate analytic series after
substituting small quantities such as \(y\) and \(y\log y\). The theorem
also does not assert that every formally well-based infinite power-log input
has a convergent inverse expansion. Finiteness of the analytic lift is an
explicit hypothesis.
\end{remark}

\subsection{Fixed derivatives and analytic remainder control}
The analytic lift is also valid in relative complex disks
\(|t-t_0|<c t_0\), with \(c\in(0,1)\) fixed and \(t_0>0\) small. A consistent
logarithm exists there, and \(|t|\) is comparable to \(t_0\) while
\(|\log t|\le |\log t_0|+C\). Cauchy estimates on such disks show that a
fixed number of derivatives of the degree tail in
\eqref{eq:degree-tail} costs the corresponding powers of \(t_0^{-1}\).
This is a reason for differentiated remainder estimates in this class.
It is \emph{not} a general license to differentiate an arbitrary real-axis
big-\(O\) statement. The distinction is important when extending the method
to tails or residuals known only on the real axis.

\section{Exact truncation and the first omitted block}
\label{sec:truncation}

\subsection{An exclusive cutoff in the output variable}
The package takes a positive real cutoff \(C\) and retains precisely the
blocks \(y^q P(\log y)\) with \(q<C\). Since the normalized block
\(t^{1+\mu}\) becomes \(y^{(1+\mu)/\alpha}\) times a constant, the relative
cutoff is
\begin{equation}\label{eq:cutoff}
 W=\alpha C-1.
\end{equation}
If \(W\le0\), even the leading term is excluded. Otherwise define
\begin{equation}\label{eq:nmax}
 N=\left\lceil\frac{W}{\delta_*}\right\rceil.
\end{equation}
No correction order \(n\ge N\) contributes below the cutoff, since its
relative valuation is at least \(n\delta_*\). Computing through order \(N\)
is sufficient to determine both the retained coefficients and a first-omitted
boundary block.

An immediate, always valid bound for the truncated inverse \(G_C\) is
\[
 g(y)-G_C(y)=\bigO\!\left(y^C(1+|\log y|)^{dN}\right).
\]
That bound is often unnecessarily weak. The implementation returns a sharper
power and a sharper logarithmic degree.

\subsection{Keeping a boundary while discarding larger powers}
At step \(n\), let \(A_n\) be the product of the retained part of \(u^{n-1}\)
with \(u\), before truncation. Split its blocks at relative weight \(W\).
The blocks below \(W\) are sufficient for the next multiplication. Among all
blocks at or above \(W\) encountered for \(1\le n\le N\), remember the
smallest weight \(\mu_0\). Apply the appropriate Euler operator to its
coefficient at every order at which it occurs, and add those contributions.
Call the resulting polynomial \(B(L)\).

This boundary is meaningful even though larger blocks are not propagated.
Any descendant of a discarded block has \emph{strictly greater} weight.
It cannot return to a smaller exponent or contribute to the smallest
omitted weight. Polynomial cancellations inside a product are performed
before recording its support.

\begin{theorem}[First-omitted block and analytic remainder]\label{thm:frontier}
Assume \(u\ne0\) and \(W>0\). The boundary procedure above terminates,
\(W\le\mu_0\le N\delta_*\), and computes the complete coefficient of the
inverse at \(t^{1+\mu_0}\). With
\[
 \rho=\frac{1+\mu_0}{\alpha},\qquad
 Q(\ell)=a^{-\rho}B\!\left(\frac{\ell-\log a}{\alpha}\right),
\]
we have
\begin{equation}\label{eq:boundary-asymptotic}
 g(y)-G_C(y)=y^\rho Q(\log y)+\smallo(y^\rho).
\end{equation}
If \(Q\ne0\), set \(k=\deg Q\); otherwise set \(k=0\). In either case,
\begin{equation}\label{eq:returned-remainder}
 g(y)-G_C(y)=\bigO\!\left(y^\rho(1+|\log y|)^k\right).
\end{equation}
A vanishing \(Q\) does not imply exactness.
\end{theorem}
\begin{proof}
The pure least-weight product in \(u^N\) has weight \(N\delta_*\) and
nonzero polynomial coefficient. Its prefixes have weight below \(W\) until
the first crossing, so some omitted block is encountered by order \(N\).
This proves the bound on \(\mu_0\).

A coefficient at \(\mu_0\) omitted from an intermediate multiplication could
only have an ancestor of weight at least \(W\). Such an ancestor, if nonzero,
is itself an encountered omitted block and has weight at least the final
minimum \(\mu_0\). Multiplying it by a positive-weight correction would
produce weight greater than \(\mu_0\), a contradiction. Thus the complete
boundary coefficient is obtained. The same argument proves that all
retained coefficients are complete.

In the full formula, the finitely many terms of degrees at most \(N\) with
weight strictly above \(\mu_0\) are \(\smallo(t^{1+\mu_0})\), since a fixed
positive power gap dominates every fixed logarithmic power. The remaining
degree tail starts at \(N+1\) and is
\[
 \bigO\!\left(t^{1+(N+1)\delta_*}(1+|L|)^{d(N+1)}\right)
 =\smallo(t^{1+\mu_0}).
\]
Theorem~\ref{thm:convergence} justifies that estimate. Converting from \(t\)
to \(y\) gives \eqref{eq:boundary-asymptotic} and then
\eqref{eq:returned-remainder}.
\end{proof}

\subsection{Why a plain \texorpdfstring{$O(y^3)$}{O(y3)} would be wrong}
For the first example, retaining only \(y-y^2(1+\log y)\) leaves
\[
 y^3(2\log^2 y+5\log y+3)+\smallo(y^3).
\]
After division by \(y^3\), this is unbounded. The correct simple bound is
\(\bigO(y^3(1+|\log y|)^2)\), not \(\bigO(y^3)\).
A formal truncation symbol based only on power valuation must not be read as
an analytic Landau estimate without its logarithmic degree.

As another example, the inverse of \(x+x^2+2x^3\) begins
\[
 y-y^2+0\cdot y^3+5y^4+\cdots.
\]
At cutoff \(3\), the boundary exponent is \(3\) but its polynomial cancels.
The package conservatively returns a bound at power \(3\), with boundary
term zero; it does not invent a proof that no later terms exist. Raising the
cutoff exposes the next nonzero coefficient.

\section{The logarithmic example in full detail}
\label{sec:log-example}

For \(f_1\), we have \(u(t)=t(1+L)\), \(a=\alpha=1\). Write
\begin{equation}\label{eq:f1-series}
 g_1(y)=y+\sum_{n\ge1}y^{n+1}P_n(\log y).
\end{equation}
Equation~\eqref{eq:block-coefficient} gives
\begin{equation}\label{eq:f1-polynomial}
 \boxed{\displaystyle
 P_n(L)=\frac{(-1)^n}{n!}
       \prod_{j=2}^{n}(n+j+\partial_L)(1+L)^n.}
\end{equation}
Equivalently, the factors are \(n+2+\partial_L,\ldots,2n+\partial_L\).
The empty product for \(n=1\) gives \(P_1=-(1+L)\).

The first coefficients, independently checked by composition, are
\begin{align}
 P_1(L)&=-1-L,\notag\\
 P_2(L)&=3+5L+2L^2,\notag\\
 P_3(L)&=-\frac{23}{2}-27L-\frac{41}{2}L^2-5L^3,\label{eq:f1-coeffs}\\
 P_4(L)&=\frac{299}{6}+151L+\frac{335}{2}L^2
                 +\frac{241}{3}L^3+14L^4,\notag\\
 P_5(L)&=-\frac{2789}{12}-\frac{5177}{6}L-1256L^2
 -\frac{5365}{6}L^3-\frac{3727}{12}L^4-42L^5.\notag
\end{align}
Thus an answer through the fourth power of \(y\) is
\begin{align}\label{eq:f1-answer}
 g_1(y)={}&y-y^2(1+L)+y^3(3+5L+2L^2)\notag\\
 &-y^4\left(\frac{23}{2}+27L+\frac{41}{2}L^2+5L^3\right)
 +\bigO\!\left(y^5(1+|L|)^4\right),\qquad L=\log y.
\end{align}

\subsection{An elementary triangular check}
Put \(g=y s\), \(s=1+c_1y+c_2y^2+\cdots\). The equation becomes
\[
 s+y s^2(1+\log y+\log s)=1.
\]
At first order, \(c_1=-(1+L)\). At second order,
\[
 c_2+(3+2L)c_1=0,
 \qquad c_2=(1+L)(3+2L).
\]
Continuing gives \eqref{eq:f1-coeffs}. This calculation shows explicitly
where the derivative of the logarithm enters: the \(\log s\) term contributes
at the next order and cannot be omitted. The independent validation script
solves this triangular system through three corrections.

\subsection{Catalan numbers in the top logarithmic degree}
Every factor in \eqref{eq:f1-polynomial} preserves degree and multiplies the
leading coefficient by its constant part. Therefore
\begin{equation}\label{eq:catalan}
 \deg P_n=n,\qquad
 [L^n]P_n(L)=(-1)^n\frac{(2n)!}{n!(n+1)!}
            =(-1)^n\,\mathrm{Cat}_n.
\end{equation}
The sequence \(-1,2,-5,14,-42,132,\ldots\) in the leading logarithms is
therefore forced by the universal formula. It also follows from keeping
only the dominant \(y\log y\) perturbation in the analytic lift.
The lower logarithmic coefficients encode the interactions with the constant
part of \(1+\log y\) and with \(\log s\).

\subsection{A two-variable convergence picture}
For this example, a particularly economical analytic lift is
\[
 s+w s^2+z s^2\log s=1,
 \qquad z=y,\quad w=y(1+\log y).
\]
The left side has derivative one with respect to \(s\) at \((s,z,w)=(1,0,0)\).
Its solution is analytic in \((z,w)\) near the origin. The path
\((y,y(1+\log y))\) tends to that origin, so the power-log expansion converges
for sufficiently small positive \(y\). This is an explicit resolution of the
apparent tension between the logarithmic endpoint singularity and the
convergence theorem.

\section{The irrational-power example and a closed coefficient law}
\label{sec:irrational}

Let \(p>1\) be any fixed real number and consider \(f_p(x)=x+x^p\).
The Euler formula gives
\begin{equation}\label{eq:fuss}
 g_p(y)=\sum_{n\ge0}b_n(p)y^{1+n(p-1)},
 \qquad
 b_n(p)=\frac{(-1)^n}{1+n(p-1)}\binom{np}{n}.
\end{equation}
For \(n\ge1\), an equivalent finite-product expression is
\begin{equation}\label{eq:fuss-product}
 b_n(p)=\frac{(-1)^n}{n!}\fall{np}{n-1}
       =\frac{(-1)^n\Gamma(np+1)}
       {\Gamma(n+1)\Gamma(n(p-1)+2)}.
\end{equation}
The finite product is convenient for exact algebra; the gamma expression is
convenient for coefficient asymptotics. The first terms are
\begin{align}\label{eq:p-general}
 g_p(y)={}&y-y^p+p\,y^{2p-1}
 -\frac{p(3p-1)}2y^{3p-2}\notag\\
 &+\frac{p(4p-1)(4p-2)}6y^{4p-3}
 -\frac{p(5p-1)(5p-2)(5p-3)}{24}y^{5p-4}+\cdots.
\end{align}
For \(p=\sqrt2\), this yields exactly the requested beginning:
\begin{align}\label{eq:sqrt2-answer}
 g_2(y)={}&y-y^{\sqrt2}+\sqrt2\,y^{2\sqrt2-1}
 -\frac{6-\sqrt2}{2}y^{3\sqrt2-2}\notag\\
 &+\left(\frac{17\sqrt2}{3}-4\right)y^{4\sqrt2-3}
 +\bigO\!\left(y^{5\sqrt2-4}\right).
\end{align}
Here the error has no logarithmic factor.

\subsection{An ordinary analytic series in a nonordinary variable}
Set \(z=y^{p-1}\) and \(g_p(y)=y H(z)\). The inverse equation is
\begin{equation}\label{eq:fuss-implicit}
 H+zH^p=1,\qquad H(0)=1.
\end{equation}
This is an ordinary analytic implicit equation near \((H,z)=(1,0)\), even
when \(p\) is irrational. Its ordinary Taylor coefficients in \(z\) are the
\(b_n(p)\). The irrational exponents appear only when substituting
\(z=y^{p-1}\) back into the result.

\subsection{Radius of convergence}
Stirling's formula applied to \eqref{eq:fuss-product} gives, with
\(\delta=p-1\),
\begin{equation}\label{eq:coeff-asymptotic}
 |b_n(p)|\sim
 \frac{\sqrt p}{\sqrt{2\pi}\,\delta^{3/2}}\,n^{-3/2}
 \left(\frac{p^p}{\delta^\delta}\right)^n.
\end{equation}
Hence the exact Taylor radius in \(z\) is
\begin{equation}\label{eq:radius}
 R_p=\frac{(p-1)^{p-1}}{p^p}.
\end{equation}
The critical point of \eqref{eq:fuss-implicit} is consistent with this:
\(H_*=p/(p-1)\) and \(z_*=-R_p\), where
\(1+p z H^{p-1}=0\).
For \(p=\sqrt2\),
\[
 R_p\approx0.4251956066518408,
 \qquad R_p^{1/(p-1)}\approx0.1268627051233097.
\]
Thus the inverse series converges, in particular, for
\(0<y<R_p^{1/(p-1)}\). The positive inverse itself exists beyond this
interval; the bound describes the Taylor series in \(z\), not the global
domain of the inverse. The \(n^{-3/2}\) factor also gives absolute convergence
on the circle \(|z|=R_p\), although the negative critical point is singular.

\section{Mixed exponents, scaling, and resonances}
\label{sec:mixed}

\subsection{A leading monomial other than \texorpdfstring{$x$}{x}}
For
\[
 f(x)=4x^2\bigl(1+x(1+\log x)\bigr),
 \qquad t=\frac{\sqrt y}{2},\quad L=\log t,
\]
the first two correction orders from \eqref{eq:master} are
\begin{equation}\label{eq:scaled-example}
 g(y)=t-\frac12t^2(1+L)
 +\frac18t^3(7+12L+5L^2)
 +\bigO\!\left(t^4(1+|L|)^3\right).
\end{equation}
The coefficient \(7+12L+5L^2\) comes from
\((5+\partial_L)(1+L)^2\). In terms of \(y\), even the first correction
contains the shift \(\log t=\tfrac12\log y-\log2\). A normalization that
keeps \(\log y\) unchanged would give wrong coefficients.

More generally, one relative correction \(x^\delta P(\log x)\) gives
\begin{align}\label{eq:one-correction-general}
 g(y)={}&t-\frac{t^{1+\delta}}\alpha P(L)\notag\\
 &+\frac{t^{1+2\delta}}{2\alpha^2}
 \left[(2\delta+1+\alpha)P(L)^2+2P(L)P'(L)\right]+\cdots.
\end{align}
This is a useful unit test for both the exponent rescaling and the logarithm
shift.

\subsection{Two generators}
For \(u(t)=c_1t^{\delta_1}+c_2t^{\delta_2}\) with \(\alpha=1\), the first
mixed correction, from \((m_1,m_2)=(1,1)\), is
\[
 (2+\delta_1+\delta_2)c_1c_2\,
 t^{1+\delta_1+\delta_2}.
\]
When the ratio of the generators is irrational, most weights are distinct;
when it is rational, multiple multiindices can meet at the same weight.
Either case is handled by the same sparse addition operation.

For example, the source
\(x+x^{\sqrt2}+x^{\sqrt3}\) has relative generators
\(\sqrt2-1\) and \(\sqrt3-1\). The package orders the resulting real
algebraic exponents exactly. The validation includes independent composition
for this genuinely two-generator example, as well as for
\(x+x^{\sqrt2}+x^2(1+\log x)\).

\subsection{Collision handling is a correctness condition}
For \(f=x+x^2+2x^3\), the contribution at inverse power \(y^3\) from the
linear perturbation order is \(-2y^3\), whereas the quadratic perturbation
order contributes \(+2y^3\). Recording the two terms under different
unevaluated exponent keys or deciding their order numerically could prevent
the cancellation. The implementation compares mathematical exponents and
adds their full coefficient polynomials before testing for zero.

\section{Residuals, source uncertainty, and conditioning}
\label{sec:residuals}

\subsection{What an independent residual check should verify}
Let \(\beta=1/\alpha\), and suppose an inverse jet has error
\(g-G=\bigO(y^\rho(1+|\log y|)^k)\), with \(G\sim g\sim t\).
The mean-value theorem and the derivative equivalent in
Theorem~\ref{thm:branch} give
\begin{equation}\label{eq:residual-scale}
 f(G(y))-y
 =\bigO\!\left(y^{\rho+1-\beta}(1+|\log y|)^k\right).
\end{equation}
Thus the symbolic checker should find no residual block with power strictly
below
\begin{equation}\label{eq:residual-cutoff}
 C_{\rm res}=\rho+1-\beta.
\end{equation}
Comparing the residual and inverse errors at the same power would be wrong
when \(\alpha\ne1\).

The package checker recomputes composition using the binomial and logarithm
identities \eqref{eq:unit-functions}. For an inner jet
\(G=c y^\beta(1+v)\), a source block is composed as
\begin{equation}\label{eq:checker-compose}
 G^q P(\log G)
 =c^q y^{\beta q}(1+v)^q
   P\bigl(\log c+\beta\log y+\log(1+v)\bigr).
\end{equation}
This algorithm does not call the Euler coefficient generator. Passing the
check is an independent algebraic consistency test. It does not by itself
prove a global branch statement, establish the assumptions supplied by the
caller, or certify a floating-point error bound.

\subsection{Transporting an asymptotic input remainder}
\begin{theorem}[Input-jet error transport]\label{thm:input-error}
Let \(f_0\) belong to \eqref{eq:class}. Suppose \(f=f_0+R\) is continuous,
strictly increasing near zero, and
\[
 R(x)=\bigO\!\left(x^{\alpha+\sigma}(1+|\log x|)^m\right),
 \qquad \sigma>0.
\]
Then the positive inverse branches satisfy
\begin{equation}\label{eq:input-error}
 f^{-1}(y)-f_0^{-1}(y)
 =\bigO\!\left(y^{(1+\sigma)/\alpha}(1+|\log y|)^m\right).
\end{equation}
\end{theorem}
\begin{proof}
Both sources have the same leading equivalent \(a x^\alpha\), so both
inverses are asymptotic to \(t=(y/a)^{1/\alpha}\). If \(g=f^{-1}(y)\) and
\(g_0=f_0^{-1}(y)\), then
\(f_0(g)-f_0(g_0)=-R(g)\). The mean-value theorem for the known smooth
source \(f_0\) divides this error by a quantity comparable to
\(t^{\alpha-1}\). The result is
\(\bigO(t^{1+\sigma}(1+|\log t|)^m)\), which is \eqref{eq:input-error}.
No derivative bound on \(R\) is needed once the stated monotonicity and
existence of the branch for \(f\) are available.
\end{proof}

For example, replacing \(\sin x\) by
\(x-x^3/6+x^5/120\) has a source error \(\bigO(x^7)\). Its inverse
coefficients below power \(y^7\) agree with those of \(\arcsin y\), but a
boundary coefficient at \(y^7\) computed from that finite polynomial is not
necessarily the boundary coefficient of the original inverse. The package
assumes that its supplied finite expression is the exact source. A caller
using an asymptotic source jet must enforce this independent error budget;
there is no implicit claim of infinite input precision.

\subsection{A posteriori numerical bounds}
For a positive approximation \(G\), if \(f'\ge m_0>0\) on an interval
containing \(G\) and the desired root, then
\begin{equation}\label{eq:posterior}
 |g(y)-G(y)|\le \frac{|f(G(y))-y|}{m_0}.
\end{equation}
For the two original examples, global slope bounds give particularly simple
versions:
\[
 |g_1-G|\le\frac{|f_1(G)-y|}{1-2e^{-5/2}},
 \qquad
 |g_p-G|\le |f_p(G)-y|\quad(p>1),
\]
provided \(G>0\). These are mathematical inequalities. A numerically
\emph{certified} evaluation additionally requires a rigorous enclosure of
the residual; ordinary high-precision arithmetic is evidence, not an
interval certificate.

The relative condition number tends to
\[
 \left|\frac{y g'(y)}{g(y)}\right|\longrightarrow\frac1\alpha.
\]
This clean normalization explains why working with \(g/t\), rather than
with a poorly scaled raw implicit equation, is helpful for both formal and
numerical algorithms.

\section{The Wolfram Language package}
\label{sec:implementation}

\subsection{Loading and the two original examples}
The implementation is in \code{Kernel/RealInverseAsymptotics.wl}. It uses no
external Wolfram packages. Load the file and keep the source and output
variables distinct and unassigned:
\begin{lstlisting}
Get["/path/to/RealInverseAsymptotics/Kernel/RealInverseAsymptotics.wl"];
Clear[x, y];

r1 = RealInverseAsymptotic[
  x + x^2 (1 + Log[x]), {x, 0}, {y, 4}];
r1["Expression"]
(* y - y^2 (1+Log[y]) + y^3 (3+5 Log[y]+2 Log[y]^2) *)
r1["Remainder"]
(* PowerLogO[y, 4, 3] *)

r2 = RealInverseAsymptotic[
  x + x^Sqrt[2], {x, 0}, {y, 4 Sqrt[2]-3}];
r2["Expression"]
(* y-y^Sqrt[2]+Sqrt[2] y^(2 Sqrt[2]-1)
   -(6-Sqrt[2])/2 y^(3 Sqrt[2]-2) *)
\end{lstlisting}
The comments show mathematically verified expected expressions, not a
transcript from an executed Wolfram kernel. This distinction applies to all
package examples in this article.

The first call retains powers strictly below \(4\), not four terms and not
all powers through \(4\) inclusively. The second cutoff is chosen to exclude
the next power \(4\sqrt2-3\), so it returns the four blocks requested in the
question. A cutoff such as \(3\) is also accepted; its meaning is the same
exclusive real-exponent bound.

\subsection{The result object}
The function returns a standard \code{Association}. Its main fields are
listed in Table~\ref{tab:fields}. Keeping the finite expression separate
from the remainder avoids pretending that a generalized series belongs to
\code{SeriesData}.

\begin{table}[htbp]
\centering\small
\begin{tabularx}{\textwidth}{@{}lX@{}}
\toprule
Field & Meaning \\
\midrule
\code{Expression} & The retained finite inverse expression. \\
\code{Terms} & Sorted pairs \(\{q,P(\log y)\}\), one per exact power. \\
\code{RemainderPower} & The first possible omitted output exponent \(\rho\). \\
\code{RemainderLogDegree} & The nonnegative integer \(k\) in the returned bound. \\
\code{RemainderScale} & \(y^\rho(1+|\log y|)^k\), or zero for an exact result. \\
\code{Remainder} & The inert annotation \code{PowerLogO[y,rho,k]}. \\
\code{BoundaryTerm} & The complete term \(y^\rho Q(\log y)\); it can cancel. \\
\code{LeadingTerm} & \((y/a)^{1/\alpha}\), the selected branch scale. \\
\code{Exact} & True only for an untruncated pure-monomial source. \\
\code{Cutoff} & The requested exclusive exponent bound. \\
\bottomrule
\end{tabularx}
\caption{Principal output fields. Source variables, source expression,
assumptions, branch information, and the finite order bound are also stored.}
\label{tab:fields}
\end{table}

\code{PowerLogO} has no arithmetic rules. It is a declaration of an
asymptotic bound, not an object that can safely be substituted, multiplied,
or differentiated without additional analysis. To evaluate numerically,
use \code{r["Expression"]} and only then apply \code{N}; do not evaluate
the exponents approximately before constructing the support.

If the cutoff excludes the leading term, the retained expression is zero
and the omitted boundary is the leading term. The residual checker refuses
such an empty inverse approximation because composition through its
logarithm would be undefined. For a pure source \(a x^\alpha\) whose leading
term is retained, the inverse is exact and the remainder is zero.

\subsection{Parser and coefficient normalization}
The parser first replaces \code{Log[x]} by a fresh polynomial coordinate,
then uses ordinary \code{Expand}. Each resulting summand must be a product
of a power \(x^q\) and an \(x\)-independent polynomial in that coordinate.
Equal exponents are merged. The smallest exponent must be a positive
constant \(\alpha\), and its coefficient must be an \(L\)-independent
constant proved positive under the supplied assumptions. Dividing the
remaining blocks by \(a x^\alpha\) constructs the positive correction gaps.

The parser deliberately does not perform branch-sensitive rewrites. For
example, normalize \(\log(2x)\) to \(\log2+\log x\) first on the positive
domain; an unreduced \code{Log[2 x]} is rejected. Noninteger powers of
products are not distributed using \code{PowerExpand}. The only power and
logarithm transformations used by the algorithm are those justified by the
explicit positive leading coefficient and the positive local branch.

All real-number literals with finite precision are rejected, even in
coefficients or the cutoff. Inputs such as \code{Sqrt[2]} and \code{7/5}
are exact; \code{1.41421356} and \code{1.4} are not. Symbolic coefficients
are allowed when their reality and the positivity of the leading
coefficient can be established, for example
\begin{lstlisting}
RealInverseAsymptotic[
  a x + a x^2, {x, 0}, {y, 3}, Assumptions -> a > 0]
\end{lstlisting}
Exponents, by contrast, must be numerically specified exact real constants.
An unspecialized symbol \(p\), even with \(p>1\), is rejected. A symbolic
exponent could change the ordering, collision pattern, and number of retained
terms as a parameter varies; a correct fully symbolic implementation would
need a parameter-space case decomposition. The present implementation does
not silently choose one such case.

\subsection{Exact exponent equality and order}
Each exponent is simplified exactly. Comparisons must resolve to less than,
equal to, or greater than; otherwise the computation returns
\code{Failure["UndecidableExponentOrder", ...]}. There is no explicit
floating-point fallback, tolerance-based equality, or machine-number hash
key. Real algebraic expressions, including the original \(\sqrt2\), are the
primary intended case. Exact transcendental constants may also work when
all necessary inequalities and iteration bounds are resolved.

The reference implementation uses sorted lists and equality-checked merging
rather than a sophisticated algebraic-number key. This is deliberately
conservative about correctness. It is not a claim that list insertion is an
optimal data structure for very large supports.

\subsection{The main loop}
The computational core implements the following finite procedure.
\begin{enumerate}
\item Parse and normalize \(f\); compute \(W=\alpha C-1\), \(\delta_*\),
and \(N=\lceil W/\delta_*\rceil\).
\item Start with the zeroth power \(u^0=1\), retained inverse coefficient
one, and no omitted boundary.
\item For \(n=1,\ldots,N\), multiply the retained part of \(u^{n-1}\) by
\(u\). For each raw block, apply \eqref{eq:block-coefficient}. Add blocks
below \(W\) to the inverse and retain their untransformed product
coefficients for the next multiplication. For omitted blocks, update the
smallest boundary weight and its summed transformed coefficient.
\item Convert \(t^{1+\mu}Q(L)\) to
\(y^{(1+\mu)/\alpha}a^{-(1+\mu)/\alpha}
 Q((\log y-\log a)/\alpha)\). Return the expression and the boundary-based
remainder data.
\end{enumerate}
The distinction between the untransformed product coefficient and the
Euler-transformed inverse coefficient is essential: the next iteration
needs \(u^n\), not the \(n\)-th inverse correction.

The operator on one block is implemented by the short loop
\begin{lstlisting}
eulerBlock[q_, mu_, alpha_, n_Integer, ell_Symbol] :=
 Module[{r = q, j},
  Do[r = poly[(mu + 1 + j alpha) r + D[r, ell]],
     {j, 1, n - 1}];
  poly[(-1)^n r/(Factorial[n] alpha^n)]
 ];
\end{lstlisting}
Here \code{poly} expands and simplifies polynomial coefficients under the
active assumptions. The surrounding package supplies parsing, merging,
cutoff handling, and failures.

\subsection{Independent composition and the checker}
Two additional public functions are supplied:
\begin{lstlisting}
PowerLogCompose[Log[x], x, y+y^2, {y,4}]
(* Expression: Log[y]+y-y^2/2+y^3/3 *)

CheckInverseAsymptotic[r1]
(* Passed -> True is the expected result. *)
\end{lstlisting}
The composition engine implements \eqref{eq:checker-compose}, using finite
binomial and logarithm sums cut off by their positive valuations. It does
not reuse the inverse coefficient formula. \code{CheckInverseAsymptotic}
uses the source expression saved in the result and the residual cutoff
\eqref{eq:residual-cutoff}. Its output includes the Boolean
\code{Passed}, the residual expression below that bound, and the bound
itself. A test deliberately corrupts an inverse coefficient to ensure that
the checker is intended to reject an incorrect jet.

\subsection{Budgets, failure behavior, and cost}
The options are
\begin{lstlisting}
Assumptions :> $Assumptions
"MaxOrder" -> 256
"MaxTerms" -> 20000
\end{lstlisting}
A temporary multiplication budget is four times \code{MaxTerms}. Exceeding
a budget returns a structured failure rather than a silently incomplete
series. These are size budgets, not a guarantee of a wall-clock running-time
bound for symbolic simplification.

If \(M\) is a bound on retained support size and the source has \(r\) blocks,
a multiplication produces at most \(rM\) raw candidates before merging.
The straightforward equality-checked merge can take quadratic time in that
candidate count. Polynomial degrees grow at most linearly with perturbation
order by \eqref{eq:degree-bound}; coefficient bit sizes and exact algebraic
comparison costs are additional factors. The implementation is appropriate
for moderate symbolic orders and for auditing the mathematics. At much
higher orders, a better support index and Newton reversion are natural
optimizations, but their complexity is not claimed for this implementation.

\section{Validation, reproducibility, and what was not executed}
\label{sec:validation}

\subsection{Independent computations that were executed}
The script \code{validation/validate.py} completed 58 checks using
SymPy~1.14.0 and mpmath~1.3.0. Its JSON report is included in the archive.
The exact checks cover the first logarithmic coefficients, the general
irrational Fuss--Catalan formula through nine correction indices, independent
composition, a separate triangular coefficient solve, exact monomial
normalization, strict cutoff edge cases, and a cancelled boundary.

Nine additional source classes were checked by independent composition:
a rational nonunit leading exponent; a nonunit leading coefficient;
resonant integer exponents; mixed logarithmic polynomials; a negative
perturbation; an irrational leading exponent; two distinct irrational
generators; an irrational/logarithmic mixture; and a polynomial logarithmic
coefficient with several signs. These tests exercise the normalization and
support logic, not only the two motivating examples.

Numerical checks used 100-digit arithmetic and positive bisection to
calculate an inverse independently of the coefficient formula. For each
original example, the script compares approximations with one, two, three,
and five correction orders at \(y=10^{-3},10^{-6},10^{-12}\).
Table~\ref{tab:numerics} shows selected values. The ratio compares the actual
absolute error with the absolute value of the computed boundary term;
its approach to one is consistent with \eqref{eq:boundary-asymptotic}.

\begin{table}[htbp]
\centering\small
\begin{tabular}{@{}llrrr@{}}
\toprule
Source & \(y\) & Orders & Absolute error & Error / boundary \\
\midrule
Logarithmic & \(10^{-3}\) & 3 & \(1.2594119\,10^{-11}\) & 1.01590491 \\
Logarithmic & \(10^{-6}\) & 3 & \(3.2814172\,10^{-25}\) & 1.00003643 \\
Logarithmic & \(10^{-12}\) & 3 & \(6.5895650\,10^{-54}\) & 1.00000000 \\
Logarithmic & \(10^{-6}\) & 5 & \(4.5579380\,10^{-34}\) & 1.00003940 \\
Irrational & \(10^{-3}\) & 3 & \(3.8876652\,10^{-8}\) & 0.90508123 \\
Irrational & \(10^{-6}\) & 3 & \(4.5691233\,10^{-16}\) & 0.99401819 \\
Irrational & \(10^{-12}\) & 3 & \(5.2638620\,10^{-32}\) & 0.99998031 \\
Irrational & \(10^{-6}\) & 5 & \(1.7139602\,10^{-20}\) & 0.99364138 \\
\bottomrule
\end{tabular}
\caption{Executed high-precision numerical checks. ``Orders'' counts
correction orders after the leading term. Values are numerical evidence,
not interval-arithmetic certificates. Full data are in
\code{validation/numerics.csv}.}
\label{tab:numerics}
\end{table}

\subsection{Wolfram tests supplied for execution}
The archive contains 35 MUnit \code{VerificationTest} cases. They check both
original examples, the logarithmic remainder, the complete omitted
polynomial, a nonunit leading monomial, an exact monomial, exponent
collisions and cancellation, symbolic positive coefficients, independent
composition, a deliberately corrupted result, and error handling.
They also specify expected rejection of inexact data, an unspecialized
symbolic exponent, a logarithmic leading scale, an essential exponential,
complex coefficients, and insufficient resource budgets.

Run the suite inside Mathematica with
\begin{lstlisting}
TestReport["/path/to/RealInverseAsymptotics/Tests/RealInverseAsymptotics.wlt"]
\end{lstlisting}
or from a command line with
\begin{lstlisting}[language={}]
wolframscript -file Tests/RunTests.wls
\end{lstlisting}

\textbf{These Wolfram tests were not executed in this session.}
The available Wolfram connector returned a tool-availability error. No
local Wolfram kernel was installed, and an attempted Mathics installation
was unavailable. Consequently there is no claimed Wolfram pass count,
no verified current-kernel compatibility claim, and no claim to have
reproduced the old \code{InverseFunction} behavior on a current release.
Independent SymPy/mpmath execution validates the derived formulas and an
independently coded algorithm; it does not execute Wolfram evaluation
semantics. A separate delimiter/comment/string balance check of the supplied
Wolfram sources is included, but it is not a Wolfram parser or runtime test.

\subsection{Reproducing the independent checks and the document}
The validation scripts and raw results are included rather than only a
summary. Their commands are
\begin{lstlisting}[language={}]
python validation/validate.py
python validation/check_wl_delimiters.py
\end{lstlisting}
The article source is self-contained apart from standard LaTeX packages.
From the \code{article} directory, compile it with
\begin{lstlisting}[language={}]
latexmk -pdf real_inverse_asymptotics.tex
\end{lstlisting}
The package and the accompanying original material in this archive are
provided under the MIT No Attribution license. The cited literature remains
under its respective terms and is not bundled as third-party full text.

\section{Extensions, alternative algorithms, and scope boundaries}
\label{sec:extensions}

\subsection{Newton reversion in the same algebra}
The normalized equation \(\Phi_t(s)=s^\alpha(1+u(ts))-1=0\) has derivative
with leading constant \(\alpha\). Newton iteration therefore has the formal
form
\begin{equation}\label{eq:newton}
 s_{k+1}=s_k-
 \frac{s_k^\alpha(1+u(ts_k))-1}
 {\alpha s_k^{\alpha-1}(1+u(ts_k))+t s_k^\alpha u'(ts_k)}.
\end{equation}
Starting from \(s_0=1\), a Taylor expansion of \(\Phi_t\) shows that an error
of positive valuation \(\eta\) is replaced by one of valuation at least
\(2\eta\). Its formal derivative is a unit, and the quadratic remainder
contains two factors of the previous error. Thus relative accuracy at least
\(2^k\delta_*\) is available after \(k\) Newton steps, when all arithmetic is
performed to the needed guard precision.

This gives an alternative algorithm, not an additional implemented backend.
It uses composition, reciprocal, and generalized unit powers; all are
available mathematically in the same algebra. The Euler formula remains
particularly useful for direct coefficients, closed forms, degree laws,
and independent checks. A future high-order implementation should state
its guard-precision rules explicitly rather than assuming that a finite
jet may be differentiated or composed without changing its error budget.

\subsection{Observables of the inverse}
The parameter identity \eqref{eq:observable-lagrange} also expands an
observable \(\Psi(g)\). For a fixed real exponent \(q\), it gives
\begin{equation}\label{eq:inverse-power}
 g(y)^q=t^q\left[1+\sum_{n\ge1}
 \frac{(-1)^n q}{n!\alpha^n}
 \prod_{j=1}^{n-1}(\Eul+q+j\alpha)u(t)^n\right].
\end{equation}
For \(q=0\) the correction vanishes. Differentiating with respect to \(q\)
at zero gives
\begin{equation}\label{eq:inverse-log}
 \log g(y)=\log t+\sum_{n\ge1}\frac{(-1)^n}{n!\alpha^n}
 \prod_{j=1}^{n-1}(\Eul+j\alpha)u(t)^n.
\end{equation}
These are useful theoretical formulas for applications that need a power,
logarithm, or derivative of the inverse rather than the inverse alone.
The shipped package obtains such finite jets through
\code{PowerLogCompose}; it does not expose separate optimized generators
for \eqref{eq:inverse-power} and \eqref{eq:inverse-log}.

\subsection{Other endpoints and branches}
A finite basepoint can be treated by introducing positive local coordinates.
For a right-hand branch through \(x_0\), write \(x=x_0+t\) and subtract the
image value \(f(x_0)\). For a left-hand branch, use \(x=x_0-t\) and account
for the sign of the output coordinate. The resulting positive-coordinate
source must then satisfy the input contract. This transformation is an
explicit caller operation; the package currently accepts the basepoint
zero only.

An asymptotic inverse at infinity can sometimes be converted to one at zero
by \(x=1/t\) and an appropriate reciprocal output variable. Whether the
transformed source lies in the finite power-log class has to be checked.
A change of variables is not a reason to discard a leading logarithmic or
exponential scale that the coefficient algebra cannot represent.

\subsection{A logarithmic leading scale requires a different normalization}
A source such as
\[
 f(x)=a x^\alpha(-\log x)^b
\]
is outside \eqref{eq:class} when \(b\ne0\). Its leading coefficient is not
constant and its inverse is not asymptotic to \((y/a)^{1/\alpha}\).
There is nevertheless a useful explicit leading-scale inversion. Put
\(\ell=-\log x\). For \(b\ne0\),
\begin{equation}\label{eq:log-leading}
 \ell=-\frac b\alpha
 W_k\!\left(-\frac\alpha b(y/a)^{1/b}\right),
 \qquad x=e^{-\ell}.
\end{equation}
For \(b>0\) and \(y\downarrow0\), the required large positive \(\ell\)
uses \(W_{-1}\); for \(b<0\), it uses \(W_0\) on a large positive argument.
This follows directly by rewriting
\((y/a)^{1/b}=\ell e^{-\alpha\ell/b}\).

For example, the small positive inverse of \(x(-\log x)\) is
\[
 x=-\frac{y}{W_{-1}(-y)}
 \sim\frac{y}{\log(1/y)}.
\]
Its corrections involve \(\log\log(1/y)\). They require a nested logarithmic
scale or an exact leading Lambert expression. The current package rejects
this input rather than misclassifying a zero correction gap as positive.
Formula~\eqref{eq:log-leading} describes a possible normalization adapter,
not an implemented universal log-leading solver.

\subsection{Exponentially small sectors and Fabius-type applications}
For \(f(x)=e^{-1/x^\beta}\), the inverse is
\((\log(1/y))^{-1/\beta}\). Its source has no positive power leading
monomial at zero. For \(f(x)=x+e^{-1/x}\), the power-log inverse jet is simply
\(y\) to every algebraic order, but the inverse is not exactly \(y\): formally,
\[
 g(y)=y-e^{-1/y}+y^{-2}e^{-2/y}+\cdots.
\]
This illustrates why an all-orders power-log jet does not determine
exponentially small corrections. The sparse package rejects the exponential
input; it does not silently replace it by its power-log projection.

For the supplied Fabius research context, an appropriate larger architecture
would first identify a leading tail scale, transform to a coordinate where
the inverse equation is near identity, then state the support and derivative
bounds of every residual retained in that coordinate. Only the resulting
power-logarithmic subproblem, when it satisfies \eqref{eq:class}, can be
handed to this package unchanged. A real-axis tail equivalent alone is not
a proof of a complex collar, of differentiated residual bounds, or of the
admissibility of a chosen transseries substitution. Those distinctions are
exactly the sort of support and realization issues emphasized in the
supplied notes~\cite{proveit}.

\subsection{The precise boundary of this implementation}
The implemented class includes finite sums of exact real powers with
polynomial logarithmic coefficients, positive monomial leading terms,
several irrational generators, and real symbolic coefficients under
sufficient assumptions. It excludes a logarithmic leading monomial,
negative powers of the logarithm, iterated logarithms, essential
exponentials, oscillatory coefficients, arbitrary implicit equations,
and unspecialized symbolic exponents. Some excluded problems are easier
than the included ones after a suitable transformation; exclusion here is
an interface and proof-contract boundary, not a mathematical impossibility.

\section{Conclusions}

The two motivating examples do not require a closed form for the inverse or
an ordinary Taylor series at zero. They require an explicitly selected
positive germ, a coefficient algebra that retains logarithmic differentiation,
and an exact real-exponent support.

For a finite power-logarithmic perturbation of \(a x^\alpha\), these needs
are met by the Euler--Lagrange formula \eqref{eq:master}. Its support is a
locally finite positive semigroup. The analytic lift proves convergence for
small positive arguments, and the first-omitted theorem turns formal
truncation into an honest logarithmic Landau bound. The sparse package
implements this finite construction and independently checks the resulting
coefficients by composition. The two original inverse expansions follow
as direct specializations, while the normalization and error-transport
results make the method usable for a substantially broader family of
real inverse problems.

\appendix
\section{Additional coefficient and implementation checks}

The sixth logarithmic polynomial, useful for checking the first omitted
term when retaining five corrections, is
\begin{align*}
 P_6(L)={}&132L^6+\frac{5981}{5}L^5+\frac{52937}{12}L^4
 +\frac{16979}{2}L^3\\
 &+9007L^2+\frac{30019}{6}L+\frac{22769}{20}.
\end{align*}
Its leading coefficient is the sixth Catalan number, consistent with
\eqref{eq:catalan}. The validation report contains this polynomial and the
raw numerical comparisons that use it as an omitted boundary term.

For a one-block source, the general formula also makes hand checks easy.
If \(u=t^\delta P(L)\), the first three relative coefficients are
\begin{align*}
 A_1&=-\frac{P}{\alpha},\\
 A_2&=\frac{1}{2\alpha^2}
       (2\delta+1+\alpha+\partial_L)P^2,\\
 A_3&=-\frac{1}{6\alpha^3}
       (3\delta+1+\alpha+\partial_L)
       (3\delta+1+2\alpha+\partial_L)P^3.
\end{align*}
They multiply \(t^{1+\delta}\), \(t^{1+2\delta}\), and
\(t^{1+3\delta}\), respectively. Setting \(\alpha=1\),
\(\delta=p-1\), and \(P=1\) recovers the irrational-power coefficients;
setting \(\delta=1\) and \(P=1+L\) recovers the logarithmic example.

\section{Archive map}
\begin{lstlisting}[language={},basicstyle=\small\ttfamily]
RealInverseAsymptotics/
  README.md
  LICENSE
  Kernel/
    RealInverseAsymptotics.wl
    init.m
  Examples/
    Examples.wl
  Tests/
    RealInverseAsymptotics.wlt
    RunTests.wls
  article/
    real_inverse_asymptotics.tex
    real_inverse_asymptotics.pdf
  validation/
    validate.py
    validation_report.json
    numerics.csv
    example_coefficients.tex
    run.log
    check_wl_delimiters.py
    wl_delimiter_report.json
\end{lstlisting}

\begin{thebibliography}{99}
\bibitem{question}
Vladimir Reshetnikov,
\emph{How to get an asymptotic of the real-valued branch of the inverse function},
Mathematica Stack Exchange, question 236367.
The supplied archive was used for the exact problem statement and historical
Wolfram evaluations.
\url{https://mathematica.stackexchange.com/questions/236367/how-to-get-an-asymptotic-of-the-real-valued-branch-of-the-inverse-function}.

\bibitem{proveit}
Vladimir Reshetnikov, \emph{ProveIt: series-and-transseries research drafts},
especially the Lagrange--B\"urmann, support, logarithmic-degree, and
precision discussions. Selected excerpts consulted 7 September 2026.
\url{https://github.com/VladimirReshetnikov/ProveIt/tree/main/Analysis/FabiusFunction/docs/semi-formalized-research-frontiers/drafts/series-and-transseries}.

\bibitem{dlmf}
NIST Digital Library of Mathematical Functions,
\S1.10(vii), \emph{Inverse Functions}, including the Lagrange inversion theorem.
\url{https://dlmf.nist.gov/1.10#vii}. Accessed 7 September 2026.

\bibitem{edgar}
G.~A. Edgar, \emph{Transseries for beginners},
Real Analysis Exchange \textbf{35}(2) (2009/2010), 407--452.
\url{https://arxiv.org/abs/0801.4877}.

\bibitem{hoeven}
Joris van der Hoeven, \emph{Transseries and Real Differential Algebra},
Lecture Notes in Mathematics 1888, Springer, 2006.
\url{https://doi.org/10.1007/3-540-35590-1}.

\bibitem{salvy}
Bruno Salvy and John Shackell,
\emph{Symbolic Asymptotics: Multiseries of Inverse Functions},
Journal of Symbolic Computation \textbf{27}(6) (1999), 543--563.
\url{https://doi.org/10.1006/jsco.1999.0281}.
Bibliographic record and abstract:
\url{https://kar.kent.ac.uk/16588/}.

\bibitem{seriesdata}
Wolfram Research, \emph{SeriesData}, Wolfram Language Documentation.
\url{https://reference.wolfram.com/language/ref/SeriesData.html}.
Accessed 7 September 2026.

\bibitem{inverseseries}
Wolfram Research, \emph{InverseSeries}, Wolfram Language Documentation.
\url{https://reference.wolfram.com/language/ref/InverseSeries.html}.
Accessed 7 September 2026.

\bibitem{asymptotic}
Wolfram Research, \emph{Asymptotic}, Wolfram Language Documentation.
\url{https://reference.wolfram.com/language/ref/Asymptotic.html}.
Accessed 7 September 2026.
\end{thebibliography}
\end{document}
